\section{The summed area table baseline}
\label{sec:baseline}

This section recaps the algorithm of \citet{faggian2015} as far as it is needed
to follow the adaptations in Section~\ref{sec:adaptations}. The notation follows
the paper.

\subsection{Summed area table and windowed sums}

A field is first reduced to a binary event field $I$ by thresholding (for the
classic FSS, $I(i,j)=1$ where the field exceeds the threshold and $0$
otherwise). The summed area table $\hat{O}$ stores at each point the sum of all
cells above and to the left, inclusive:
%
\begin{equation}
  \hat{O}(i,j) = \sum_{\hat{\imath}\le i}\;\sum_{\hat{\jmath}\le j}
                 O(\hat{\imath},\hat{\jmath}).
  \label{eq:sat}
\end{equation}
%
In the code this is a double cumulative sum (\codeword{compute_integral_table}):

{\setstretch{1}\begin{lstlisting}[style=python]
def compute_integral_table(field):
    if field.ndim == 2:
        return field.cumsum(1).cumsum(0)
    elif field.ndim == 3:
        return field.cumsum(2).cumsum(1)
\end{lstlisting}}

Once $\hat{O}$ is available, the sum over a window of half-width $w$ centred on
$(i,j)$ is obtained from just four lookups,
%
\begin{equation}
  \sum_{\hat{\imath}=i-w}^{i+w}\;\sum_{\hat{\jmath}=j-w}^{j+w}
  O(\hat{\imath},\hat{\jmath})
  = \hat{O}(i+w, j+w) + \hat{O}(i-w, j-w)
  - \hat{O}(i-w, j+w) - \hat{O}(i+w, j-w),
  \label{eq:boxsum}
\end{equation}
%
\emph{independent of the window size}. This is the key property that makes the
SAT method cheap when many window sizes are evaluated per threshold: the table
is built once, and every window size is then four array lookups per grid point.
In the code the windowed sum is \codeword{integral_filter}.

\subsection{Boundary conditions}

Windows centred near the edge of the domain extend beyond the indexable region
of the table. \citet{faggian2015} clip the corner coordinates of
Eq.~\eqref{eq:boxsum} to the valid index range, which is how the reference
implementation (and \texttt{fss\_SAT.py}) handle the boundary. The reference
code realises the clip with \codeword{np.mgrid} and \codeword{np.clip}:

{\setstretch{1}\begin{lstlisting}[style=python]
r, c = np.mgrid[0:rows, 0:cols]
r0, c0 = (np.clip(r - w, 0, rows - 1), np.clip(c - w, 0, cols - 1))
r1, c1 = (np.clip(r + w, 0, rows - 1), np.clip(c + w, 0, cols - 1))
integral_table  = field[..., r1, c1] + field[..., r0, c0]
integral_table -= field[..., r0, c1] + field[..., r1, c0]
\end{lstlisting}}

\subsection{From windowed sums to the score}

Dividing each windowed sum by the window area $(2w+1)^2$ gives the fractions
$p_f$ and $p_o$ of Eq.~\eqref{eq:fss}. The reference \codeword{fss} routine
forms the fractions explicitly and evaluates the numerator and denominator with
\codeword{np.nanmean}:

{\setstretch{1}\begin{lstlisting}[style=python]
inv_window_area = 1. / (2. * w + 1.) ** 2
fhat = fhat * inv_window_area
ohat = ohat * inv_window_area
num   = np.nanmean(np.power(fhat - ohat, 2))
denom = np.nanmean(np.power(fhat, 2) + np.power(ohat, 2))
ret = num, denom, 1. - num / denom
\end{lstlisting}}

The higher-level routines (\codeword{fss_threshold}, \codeword{fss_cumsum_*},
\codeword{CWFSS}) build the binary tables once per threshold and call
\codeword{fss}/\codeword{integral_filter} for every requested window. This is
the structure that Section~\ref{sec:adaptations} modifies.
